\documentclass[11pt, a4paper]{article}

% ─── Packages ────────────────────────────────────────────────────────────────
\usepackage[margin=1in]{geometry}
\usepackage{amsmath, amssymb}
\usepackage{booktabs}
\usepackage{graphicx}
\usepackage{array}
\usepackage{caption}
\usepackage{hyperref}
\usepackage{parskip}
\usepackage{titlesec}
\usepackage[T1]{fontenc}
\usepackage{lmodern}

\hypersetup{
    colorlinks=true,
    linkcolor=black,
    urlcolor=blue,
    citecolor=black
}

% ─── Title ───────────────────────────────────────────────────────────────────
\title{\textbf{Sentiment Analysis on Amazon Product Reviews}\\
       \large LSTM Baseline --- CS6120 NLP Final Project (Team 102)}
\author{Xuyang Shen}
\date{March 23, 2026}

% ─── Document ────────────────────────────────────────────────────────────────
\begin{document}

\maketitle
\tableofcontents
\vspace{1em}
\hrule
\vspace{1em}

% ─────────────────────────────────────────────────────────────────────────────
\section{Data Preprocessing \& Embeddings}
% ─────────────────────────────────────────────────────────────────────────────

The raw Amazon Software reviews were preprocessed via the following pipeline:

\begin{enumerate}
    \item \textbf{Lowercasing \& punctuation removal:} All text was converted to lowercase and stripped of punctuation using Python's \texttt{string.punctuation}.
    \item \textbf{Stopword filtering:} NLTK's English stopword corpus was used to remove non-informative tokens, reducing noise and improving feature saliency.
    \item \textbf{Whitespace tokenization:} Remaining tokens were split on whitespace.
\end{enumerate}

Star ratings were converted to binary sentiment labels: ratings of 4--5 were labelled \textbf{Positive} (1), ratings of 1--2 were labelled \textbf{Negative} (0), and 3-star reviews were excluded. The dataset was split into Train / Validation / Test at an 80 / 10 / 10 ratio, yielding 8,964 / 1,121 / 1,121 samples.

\subsection*{GloVe Embeddings}

Pre-trained GloVe embeddings (300-dimensional, 6B corpus) were used to build a fixed embedding matrix. A vocabulary of \textbf{16,199 words} was constructed from the training set (minimum frequency $\geq 2$), and each word was mapped to its corresponding GloVe vector. Out-of-vocabulary (OOV) tokens were initialized with a random normal distribution ($\sigma = 0.6$).

% ─────────────────────────────────────────────────────────────────────────────
\section{Model Architecture}
% ─────────────────────────────────────────────────────────────────────────────

The model follows a standard sequence classification architecture:

\begin{table}[h]
\centering
\caption{LSTM Model Architecture}
\begin{tabular}{lll}
\toprule
\textbf{Layer} & \textbf{Configuration} & \textbf{Notes} \\
\midrule
Embedding       & 16,199 $\times$ 300 & GloVe weights, \textbf{frozen} (\texttt{requires\_grad=False}) \\
LSTM (2 layers) & hidden dim = 256    & \texttt{batch\_first=True}; dropout = 0.5 between layers \\
Dropout         & $p = 0.5$           & Applied to final hidden state \\
Linear          & $256 \rightarrow 1$ & Fully connected classification head \\
Sigmoid         & ---                 & Outputs probability $\in (0, 1)$ \\
\bottomrule
\end{tabular}
\end{table}

The embedding layer is frozen to prevent overfitting on the 10K-sample subset. The final hidden state of the last LSTM layer ($\mathbf{h}_T^{(2)}$) is passed through a dropout layer and then a linear layer, yielding a scalar probability for binary classification.

% ─────────────────────────────────────────────────────────────────────────────
\section{Training Process}
% ─────────────────────────────────────────────────────────────────────────────

\begin{table}[h]
\centering
\caption{Training Hyperparameters}
\begin{tabular}{ll}
\toprule
\textbf{Hyperparameter} & \textbf{Value} \\
\midrule
Optimizer        & Adam \\
Learning rate    & $1 \times 10^{-3}$ \\
Loss function    & Binary Cross Entropy (\texttt{BCELoss}) \\
Batch size       & 64 \\
Max epochs       & 10 \\
Early stopping patience & 3 (monitoring validation loss) \\
Hardware         & Apple MPS (Mac GPU acceleration) \\
\bottomrule
\end{tabular}
\end{table}

Early stopping monitored validation loss with a patience of 3. The best validation loss ($\approx 0.497$) was achieved at \textbf{Epoch 4}, and training was halted at \textbf{Epoch 7} as performance continued to deteriorate. The optimal model weights were restored prior to evaluation.

\begin{figure}[h]
    \centering
    \includegraphics[width=0.7\textwidth]{loss_curves.png}
    \caption{Training and Validation Loss over Epochs}
    \label{fig:loss}
\end{figure}

% ─────────────────────────────────────────────────────────────────────────────
\section{Evaluation and Results}
% ─────────────────────────────────────────────────────────────────────────────

The model was evaluated on the 10\% held-out test set (1,121 reviews).

\begin{table}[h]
\centering
\caption{Test Set Performance Metrics}
\begin{tabular}{ll}
\toprule
\textbf{Metric} & \textbf{Value} \\
\midrule
Accuracy  & 80.37\% \\
Precision & 80.50\% \\
Recall    & 99.67\% \\
F1-Score  & 0.8907  \\
\bottomrule
\end{tabular}
\end{table}

\begin{table}[h]
\centering
\caption{Confusion Matrix (Predicted vs.\ Actual)}
\begin{tabular}{cc|cc}
\toprule
 & & \multicolumn{2}{c}{\textbf{Predicted}} \\
 & & Negative (0) & Positive (1) \\
\midrule
\multirow{2}{*}{\textbf{Actual}}
  & Negative (0) & \textbf{5}   & 217 \\
  & Positive (1) & 3            & \textbf{896} \\
\bottomrule
\end{tabular}
\end{table}

\subsection{Class Imbalance Analysis}

The confusion matrix reveals a severe class imbalance in the test set. Of the 1,121 samples, 899 (80.2\%) are positive and only 222 (19.8\%) are negative. This imbalance propagates from the raw Amazon review distribution, where high-star ratings dominate in the Software category.

As a result, the model learned a heavily positive-biased decision boundary:
\begin{itemize}
    \item \textbf{Positive class recall:} 896 / 899 = 99.67\%
    \item \textbf{Negative class recall:} 5 / 222 = 2.25\%
\end{itemize}

The model classifies nearly all inputs as positive, effectively failing to detect negative sentiment.

\subsection{Discussion}

The 80.37\% accuracy appears acceptable at first glance, but is largely explained by the class prior: a trivial classifier that always predicts positive would achieve $\approx$80.2\% accuracy. The more informative metric is the near-zero true negative rate, which indicates the model has not meaningfully learned to distinguish negative sentiment.

This performance gap relative to the expected $\sim$85\% can be attributed to two factors:

\begin{enumerate}
    \item \textbf{Unaddressed class imbalance:} Standard \texttt{BCELoss} weights all samples equally, so the model minimizes loss by favoring the majority (positive) class.
    \item \textbf{Static embeddings:} Frozen GloVe vectors cannot adapt to the nuanced, often sarcastic language patterns prevalent in negative software reviews.
\end{enumerate}

A natural mitigation would be to apply a weighted loss function via the \texttt{pos\_weight} parameter in \texttt{BCEWithLogitsLoss}, penalizing false negatives more heavily. This limitation provides a clear motivation for evaluating BERT in the companion section: BERT's dynamic contextual attention mechanism may better handle the linguistic complexity of minority-class reviews.

\end{document}
